 WebSocket \cite{WebSocket} to protokół komunikacyjny, który umożliwia dwukierunkową komunikację w czasie rzeczywistym pomiędzy serwerem, a klientem. Ów protokół jest jest szczególnie przydatny do komunikacji w aplikacjach sieciowych, które potrzebują dynamiczną wymianę danych w czasie rzeczywistym.
\section{Działanie WebSocket}
Protokół WebSocket rozpoczyna swoją pracę od \textit{handshake HTTP}. Klient inicjuje połączenie z WebSocketem wysyłając specjalne żądanie HTTP do serwera. Serwer otrzymuje to żądanie, a następnie decyduje czy chce je zaakceptować, następnie wysyła odpowiedź HTTP. Jeżeli zaakceptował żądanie to przy odpowiedzi ustawia trwałe połączenie. Po nawiązaniu połączenia obydwie strony mogą przesyłać dane bez potrzebny ponownego połączenia.

\section{Zalety WebSocket}
WebSocket posiada wiele zalety w stosunku do innych tradycyjnych protokołów komunikacyjnych. Porównano zalety nad protokołem HTTP:
\begin{enumerate}
    \item \textbf{Dwukierunkowa komunikacja w czasie rzeczywistym:} dane mogą być przesyłane w obu kierunkach. Protokół HTTP umożliwia tylko zapytanie do serwera i odpowiedź z jego strony. Jest to przydatne w przypadku asynchronicznych zdarzeń wymagających komunikacji. Dodatkowo redukuje to opóźnienia.
    \item \textbf{Stałe połączenie:} po nawiązaniu połączenia nie trzeba ponownie tego robić, co znacząco zmniejsza zmniejsza opóźnienia i zużycie zasobów.
    \item \textbf{Mniejsze zużycie pasma:} WebSocket pozwala na przesyłanie czystych danych bez nagłówków HTTP, co znacząco redukuje ilość przesyłanych danych.
\end{enumerate}